We compare the terrain of Connecting Ridge, a candidate Artemis landing region at the lunar south
pole, and of two Martian sites, Jezero crater and Arcadia Planitia, with Earth surfaces. A deterministic fingerprint of relief
and texture, the structure function at hectometre lags and its Hurst exponent, ranks
\ComparedCount{} Earth windows from two declared regions for each target. The closest Earth window
to the lunar target lies at a relief distance of \BestReliefDistance{} in root mean square log
units, farther from the target than any other window of the lunar site is, and an inspection shows it to be a nearest
neighbour in this metric rather than a morphological analog. The Martian shortlists lie at relief
distances from \JezeroShortReliefMin{} to \JezeroShortReliefMax{} for Jezero and
\ArcadiaShortReliefMin{} to \ArcadiaShortReliefMax{} for Arcadia, below the median internal distance of those sites. No radar pair was qualified, so the radar comparison was not completed. A terrain encoder with a physics target and gradient reversal against the body, trained at a reduced scope, fails a promotion rule fixed before training (\GateFiveStatus): an untrained copy of the network matches its physics probe. It stays experimental and outside the ranking.
All numbers are generated from committed artifacts.
